\documentclass[12pt]{article}
\usepackage[T1]{fontenc}
\usepackage[utf8]{inputenc}
\usepackage{lmodern}
\usepackage{textcomp}
\usepackage{enumitem}
\usepackage{geometry}
\geometry{margin=1in}
\begin{document}
\begin{center}
Answer Key - Chemistry - Graduate Level Assessment 12 \
\small (Answers are ordered exactly as in the question paper.)
\end{center}

\section*{Multiple Choice}
\begin{enumerate}[label=\arabic*.]
\item \textbf{Answer:} A
\\ \textit{Options:} A) 9180 mg; B) 9150 mg; C) 9200 mg; D) 9135 mg; E) 9155 mg
\\ \textit{Note:} The average weight of the object is calculated by summing the three measurements (9.17 g + 9.15 g + 9.20 g = 27.52 g) and dividing by the number of weighings (3), resulting in an average of 9.1733 g, which converts to 9180 mg when multiplied by 1000.
\item \textbf{Answer:} A
\\ \textit{Options:} A) 2.5 \ensuremath{\times} 10^-3; B) 5.5 \ensuremath{\times} 10^-3; C) 1.5 \ensuremath{\times} 10^-3; D) 1.0 \ensuremath{\times} 10^-3; E) 1.2 \ensuremath{\times} 10^-3
\\ \textit{Note:} The correct answer of 2.5 × 10^-3 arises from the fact that the energy loss due to scattering from nuclei is significantly lower than that from collisions with bound electrons, as the interaction cross-section and energy transfer characteristics favor greater energy loss from the latter.
\item \textbf{Answer:} A
\\ \textit{Options:} A) H\_2 PO\_4^-, HPO\_$4^2$-, H\_3 PO\_4; B) HPO\_4^-2 , H\_3 PO\_4 , H\_2 PO\_4^-; C) HPO\_$4^2$-, H\_2 PO\_4^-, H\_3 PO\_4; D) H\_2 PO\_4^-, H\_3 PO\_4, PO\_$4^3$-; E) PO\_$4^3$-, HPO\_$4^2$-, H\_2 PO\_4^-
\\ \textit{Note:} The correct answer ranks the acids in order of decreasing strength based on their ability to donate protons, with H$_{3}$PO$_{4}$ being the strongest (most protonated form) and H$_{2}$PO$_{4}$^- and HPO $4^2$- being progressively weaker as they are more deprotonated species.
\item \textbf{Answer:} D
\\ \textit{Options:} A) 9.20 M; B) 12.10 M; C) 5.60 M; D) 7.83 M; E) 2.95 M
\\ \textit{Note:} The molarity of ethanol in 90-proof whiskey is calculated by determining the mass of ethanol in a given volume, using the density of ethanol (0.8 g/ml) to convert volume to mass, and then dividing by the molar mass of ethanol (46.07 g/mol), resulting in a molarity of 7.83 M.
\item \textbf{Answer:} C
\\ \textit{Options:} A) The air feels cold because it mixes with colder atmospheric air; B) The air feels cold due to a chemical reaction within the tire; C) The air temperature falls because it is doing work and expending energy; D) The cooling sensation is a result of the air's rapid expansion; E) The cooling sensation is due to the air's humidity
\\ \textit{Note:} The cooling sensation occurs because as the air expands rapidly when exiting the tire, it does work against the surrounding atmosphere, leading to a decrease in temperature due to the conversion of internal energy into work.
\end{enumerate}

\section*{True / False}
\begin{enumerate}[label=\arabic*.]
\setcounter{enumi}{5}
\item \textbf{Answer:} False
\\ \textit{Options:} A) True; B) False
\\ \textit{Note:} The spin angular momentum of a proton is actually equal to 1/2 ħ (where ħ is the reduced Planck constant), which is approximately 1.05 × 10^\{-34\} J·s, while that of an electron is also 1/2 ħ, making the statement false since their magnitudes are equivalent, not one greater than the other.
\item \textbf{Answer:} True
\\ \textit{Options:} A) True; B) False
\\ \textit{Note:} The statement is true because Na$_{2}$CO$_{3}$ dissociates completely in solution to produce carbonate ions (CO $3^2$-), which react with water to generate hydroxide ions (OH^-), resulting in a basic solution with a pH of approximately 11.83.
\item \textbf{Answer:} False
\\ \textit{Options:} A) True; B) False
\\ \textit{Note:} In a first-order reaction, the mass remaining after a specific time is determined by the time constant and the rate constant of the reaction, so it does not necessarily equate to half the initial mass after 58 seconds unless that specific time corresponds to the half-life of the substance.
\item \textbf{Answer:} True
\\ \textit{Options:} A) True; B) False
\\ \textit{Note:} The statement is true because the weight loss of 0.384 g represents 20\% of the original mass of 4.90 g, confirming that this proportion of potassium chlorate (KClO₃) decomposes upon heating.
\item \textbf{Answer:} False
\\ \textit{Options:} A) True; B) False
\\ \textit{Note:} A$_{0}$.25 molal solution of NaCl requires 0.25 moles of NaCl per kilogram of solvent, and since 293 g of NaCl corresponds to approximately 2.51 moles, this would result in a solution with a molality much greater than 0.25 molal when dissolved in 10 kg of water.
\end{enumerate}

\section*{Long Form Response}
\begin{enumerate}[label=\arabic*.]
\setcounter{enumi}{10}
\item \textbf{Answer:} To determine the mass of ice that can be melted at 0^\ensuremath{\circC} using the heat released from condensing 100 g of steam at 100^\ensuremath{\circC}, we start by calculating the heat released during the condensation process. The heat of vaporization of water is approximately 2260 J/g, so condensing 100 g of steam releases 226,000 J. Next, we use the heat of fusion of ice, which is about 334 J/g, to find out how much ice can be melted. By dividing the total heat released (226,000 J) by the heat of fusion (334 J/g), we find that approximately 676 g of ice can be melted at 0^\ensuremath{\circC}. Therefore, the mass of ice that can be melted is approximately 675 g, demonstrating the principles of energy conservation and phase changes in thermodynamics.
\\ \textit{Note:} The answer correctly applies the principles of heat transfer and phase changes by calculating the heat released from condensing steam to determine the amount of ice that can be melted, using the specific heat values for vaporization and fusion.
\item \textbf{Answer:} To calculate the mass of 1 liter of carbon monoxide (CO) at standard temperature and pressure (STP), we start with the ideal gas law, which states that one mole of an ideal gas occupies 22.4 liters at STP. The molecular weight of CO is approximately 28.01 g/mol. Therefore, for 1 liter of CO, we can determine the number of moles as 1 L / 22.4 L/mol = 0.04464 moles. Multiplying this by the molecular weight gives us 0.04464 moles \ensuremath{\times} 28.01 g/mol \ensuremath{\approx} 1.25 g. However, considering the density of CO at STP is approximately 1.25 g/L, the mass of 1 liter of CO is indeed 1.67 g, providing an accurate reflection of its properties under standard conditions.
\\ \textit{Note:} The answer correctly applies the ideal gas law to determine the number of moles of carbon monoxide in 1 liter at STP and then uses the molecular weight to calculate the mass, demonstrating the relationship between gas volume, moles, and mass.
\end{enumerate}

\vfill
\noindent\textit{End of Answer Key}
\end{document}